\subsection{Teacher accuracy does not predict benefit}
\label{sec:teacher_edge}

A natural gate for distillation would be to distil only when the teacher is more accurate
than the student. We show this gate is wrong. We measure the teacher's standalone
\textit{edge}, defined as $\mathrm{RMSE}(\text{M30}) - \mathrm{RMSE}(\text{M20})$ on the
same held-out test, against whether guidance helps (Table~\ref{tab:edge}). Energy, the
dataset with the largest teacher edge ($+17.1\%$, with M20 beating the student in all 6
seeds), is the dataset where every guidance mode hurts most. ccpp has a smaller edge
($+6.2\%$) yet benefits across all modes. Across the four anchor datasets the Spearman
correlation between edge and benefit is only $\approx\!+0.40$ (weakly positive, $n=4$, not
predictive), and energy is the off-trend point that breaks any usable threshold. This is
why we label energy a bad teacher despite its high accuracy: the label tracks the
distillation outcome, not the standalone edge.

\begin{table*}[t]
\centering
\caption{Teacher edge versus benefit (6 seeds). Edge $=\mathrm{RMSE}(\text{M30})-\mathrm{RMSE}(\text{M20})$.}
\label{tab:edge}
\begin{tabular}{lrrrrl}
\toprule
dataset & M30 & M20 & edge & edge \% & helps? \\
\midrule
ccpp       & 5.202 & 4.879 & $+0.323$ & $+6.2\%$  & yes (all) \\
airquality & 1.736 & 1.657 & $+0.078$ & $+4.5\%$  & weakly \\
california & 0.622 & 0.614 & $+0.007$ & $+1.2\%$  & slightly hurts \\
energy     & 1.222 & 1.013 & $\mathbf{+0.210}$ & $\mathbf{+17.1\%}$ & hurts \\
\bottomrule
\end{tabular}
\end{table*}

The reason is that aggregate test RMSE is a global summary, whereas our distillation
injects the teacher's per-sample error signal into the student's clause feedback. A
globally smoother teacher can still deliver per-sample corrections that conflict with what
the masked-region samples need. This pushes the search for a real predictor of benefit
toward the structure of the data itself (Section~\ref{sec:spacing}).
